\bookchapter{Capitolo 17: Interazione del telefono con altri dispositivi}

\chapterintro{Il telefono che hai in mano è in realtà un super telecomando e un centro dati. Non solo fa telefonate, ma può collegare computer, TV, orologi, altoparlanti e automobili, così che questi dispositivi non siano più isolati. In questo capitolo parliamo di come farli “fare squadra”.}

\booksection{Collaborazione con il computer: telefono e computer non sono più due mondi}

A dire il vero, la cosa che mi dava più fastidio era scattare una foto col telefono e doverla trasferire al computer per ritoccarla: serviva il cavo dati, bisognava aspettare il trasferimento. Poi ho scoperto che in realtà il cavo non serve affatto: una buona integrazione è dieci volte più veloce del cavo.

\begin{itemize}
  \item \textbf{Ecosistema Apple (AirDrop + Appunti universali)}: se usi iPhone con Mac, AirDrop è la prima funzione che devi imparare. Sul telefono tocchi Condividi AirDrop scegli il tuo Mac e la foto passa istantaneamente, senza consumare traffico o Wi-Fi, usando Bluetooth e Wi-Fi Direct. C'è anche una funzione spesso trascurata chiamata "Appunti universali": copi un testo sul telefono e sul Mac lo incolli direttamente, sembra magia.
  \item \textbf{Android + Windows (Link a Windows)}: Microsoft ha un'app chiamata "Link a Windows", integrata nei telefoni Samsung, mentre per gli altri Android basta scaricarla dallo store. Dopo l'abbinamento puoi vedere le notifiche del telefono direttamente sul computer, rispondere agli SMS, sfogliare le foto e persino fare e ricevere chiamate. La uso da due anni: in riunione non devo continuare a guardare il telefono, basta un'avviso sullo schermo del PC.
  \item \textbf{Collaborazione multi-schermo di marca}: la "Collaborazione multi-schermo" di Huawei e il "Mi Smart Hub" di Xiaomi sono ancora più spinte: proiettano lo schermo del telefono direttamente sul computer. Puoi trascinare file dal telefono al computer col mouse, e anche scorrere Douyin o rispondere a WeChat dal computer. Copia-incolla tra dispositivi e trascinamento file sono possibili, con un'esperienza davvero fluida.
\end{itemize}

\booksection{Proiezione sul TV: per guardare tutti insieme}

Durante le feste a casa, mia cugina voleva riprodurre sul TV il video del suo viaggio di laurea sul telefono. Dopo aver armeggiato a lungo, abbiamo scoperto che è davvero semplicissimo:

\begin{itemize}
  \item \textbf{iPhone TV}: scorri verso il basso dal lato in alto a destra per aprire il Centro di Controllo, tocca l'icona con i due riquadri sovrapposti (Screen Mirroring), trova il nome del tuo TV e tocca per collegarti. Il presupposto è che il TV e il telefono siano sulla stessa rete Wi-Fi e che il TV supporti AirPlay (la maggior parte dei TV smart lo supporta).
  \item \textbf{Android TV}: ogni marchio lo chiama in modo diverso, Huawei lo chiama "Proiezione wireless", Xiaomi "Cast", OPPO/vivo "Interazione multi-schermo". In genere si trova nel pannello degli interruttori rapidi; se non lo trovi, cerca "proiezione" nelle Impostazioni.
  \item \textbf{Proiezione diretta dall'app video}: se vuoi solo guardare iQIYI, Bilibili o YouTube sul TV, è ancora più semplice: nell'interfaccia di riproduzione cerca una piccola icona a forma di rettangolo con il simbolo Wi-Fi, toccala, scegli il TV e il video passa al TV. Il telefono nel frattempo può fare altro, il video continua sul TV senza interferenze.
  \item \textbf{Cosa fare se la proiezione non riesce}: nel 90\% dei casi non sei sulla stessa rete Wi-Fi. Nel restante 10\% la funzione di proiezione del TV non è attiva: cerca l'interruttore "Interazione multi-schermo" o "AirPlay" nelle impostazioni del TV.
\end{itemize}

\booksection{Con orologi smart e braccialetti: il tutor di salute al polso}

Ho regalato un braccialetto a mia madre. All'inizio diceva: “Non faccio sport, a cosa mi serve?”. Dopo due settimane ha scoperto che dormiva solo cinque ore e ha iniziato a correggere autonomamente il proprio ritmo. Questo è il valore più concreto di orologi e braccialetti: trasformare dati di salute invisibili in numeri visibili.

\begin{itemize}
  \item \textbf{Filtro notifiche}: quando il telefono vibra senza sosta in borsa o in tasca, alza il polso per dare un'occhiata: è un WeChat del capo o una promozione di Pinduoduo? Lo capisci in un secondo; i messaggi che non meritano di tirare fuori il telefono possono aspettare.
  \item \textbf{Sincronizzazione dati salute}: l'orologio monitora 24 ore su 24 frequenza cardiaca, ossigenazione, sonno e passi; questi dati vengono riassunti automaticamente nell'app Salute sul telefono. Col tempo vedi le tue tendenze: questo mese ho camminato troppo poco? Questa settimana la qualità del sonno è migliorata?
  \item \textbf{Sincronizzazione tessera trasporti/tessera di accesso}: la tessera dell'autobus e quella di accesso configurate sul telefono, nella maggior parte dei marchi, vengono sincronizzate automaticamente sull'orologio. In metropolitana non devi tirare fuori il telefono, basta avvicinare il polso al tornello: sembra davvero che "il futuro sia arrivato".
  \item \textbf{Trova-telefono}: gli orologi hanno quasi tutti la funzione "trova telefono": lo tocchi e il telefono suona. La uso molto più del previsto: appoggio il telefono distrattamente e poi non ricordo dov'è.
\end{itemize}

\booksection{Con la casa intelligente: il telefono è l'“interruttore generale” di casa}

La casa intelligente sembra chissà cosa, ma in fondo sono solo tre cose: controllo remoto tramite telefono, regole di automazione e comando vocale.

Ti racconto alcuni miei usi concreti, così capisci quanto sia pratico:

\begin{itemize}
  \item \textbf{Climatizzatore}: sulla via del ritorno in estate, accendo in anticipo il climatizzatore dal telefono, così quando entro in casa è già fresca. Spendendo un centinaio di yuan per una presa intelligente per climatizzatore si realizza la cosa, senza dover cambiare il climatizzatore.
  \item \textbf{Luci}: a letto prima di dormire, se non hai voglia di alzarti per spegnere, tocchi sul telefono o dici "Hey Siri / Xiao Ai, spegni la luce".
  \item \textbf{Robot aspirapolvere}: prima di uscire tocchi "Avvia pulizia" sul telefono e torni con i pavimenti puliti. Puoi anche impostare una pulizia giornaliera programmata.
  \item \textbf{Telecamera}: se in casa ci sono anziani, bambini o animali, dal telefono puoi dare un'occhiata alla situazione di casa in qualsiasi momento.
\end{itemize}

Quale ecosistema scegliere? In sintesi:

\begin{itemize}
  \item Se usi iPhone preferisci dispositivi che supportino Apple HomeKit (più costosi ma con migliore protezione della privacy).
  \item Se hai molti dispositivi e cerchi rapporto qualità-prezzo l'ecosistema Xiaomi Mi Jia è il più completo: luci, prese, robot aspirapolvere, di tutto.
  \item Se usi un telefono Huawei preferisci HarmonyOS Connect, l'esperienza con i dispositivi Huawei è la migliore.
\end{itemize}

\booksection{Cuffie e altoparlanti Bluetooth: la libertà senza fili}

Passare da cuffie con cavo a cuffie Bluetooth è stato forse il miglioramento più evidente nell'uso del telefono. Non tiro più fuori un groviglio di cavi dalla borsa e non rischio che il cavo si impigli nelle maniglie delle porte.

\begin{itemize}
  \item \textbf{L'abbinamento è semplice}: la prima volta che accendi le cuffie, in genere entrano automaticamente in modalità abbinamento (lampeggio), il telefono con il Bluetooth attivo le rileva subito, tocchi e sei collegato. Da quel momento ogni volta che le tiri fuori si collegano automaticamente, non devi fare nulla.
  \item \textbf{La stessa marca ha vantaggi}: gli AirPods con iPhone mostrano automaticamente un popup con il livello di batteria e possono passare senza interruzioni tra più dispositivi Apple. Anche le Huawei FreeBuds con telefono Huawei e le cuffie Xiaomi con telefono Xiaomi hanno popup e ottimizzazioni simili.
  \item \textbf{Altoparlante Bluetooth}: per feste in famiglia o campeggio all'aperto, un altoparlante Bluetooth raddoppia l'atmosfera. Una volta collegato, il telefono diventa la tua consolle DJ.
  \item \textbf{Attenzione alla latenza}: quando giochi con le cuffie Bluetooth potrebbe esserci una leggera latenza (senti lo sparo mezzo secondo dopo). Se sei sensibile alla latenza, usa cuffie con cavo oppure compra cuffie Bluetooth con "modalità a bassa latenza".
\end{itemize}

\begin{tcbinfo}{Tecnologia UWB (Ultra-Wideband): una "percezione spaziale" più precisa del Bluetooth}
Oltre al Bluetooth, nei nuovi telefoni si nasconde un chip UWB (Ultra-Wideband). La sua capacità è la\textbf{misurazione precisa delle distanze e il posizionamento spaziale}: riesce a calcolare quanto lontano e in quale direzione si trova un altro dispositivo, con precisione fino al centimetro. Il Bluetooth ti dice solo “è vicino”, UWB ti dice “è a 1,2 m davanti a te sulla sinistra”. Applicazioni comuni:

\begin{itemize}
  \item \textbf{Ricerca precisa di AirTag}: quando l'iPhone cerca un AirTag mostra la navigazione a freccia "7 metri, vai a sinistra", grazie a UWB. Il Bluetooth normale sa solo dire "è o non è nel raggio".
  \item \textbf{Chiave digitale dell'auto}: iPhone e flagship Android con UWB possono fare da chiave dell'auto: avvicinandosi si sblocca, allontanandosi si blocca, e riesce a distinguere il posto del guidatore da quello del passeggero, prevenendo attacchi di tipo relay.
  \item \textbf{Navigazione indoor}: nei centri commerciali e negli aeroporti, tramite beacon UWB si ottiene una navigazione "a livello di corridoio", molto più precisa del GPS.
  \item \textbf{Modelli supportati}: iPhone 11 e successivi montano tutti il chip UWB (U1/U2); in campo Android alcuni flagship di Xiaomi, Samsung, vivo e OPPO lo hanno già.
\end{itemize}

\end{tcbinfo}

\begin{tcbinfo}{LE Audio e Auracast: la nuova generazione dell'audio Bluetooth}
Il classico audio Bluetooth (A2DP), usato per vent'anni, sta per essere sostituito da\textbf{LE Audio}(audio a basso consumo), il nuovo standard audio introdotto con il Bluetooth 5.2:

\begin{itemize}
  \item \textbf{Più risparmio energetico e qualità audio migliore}: adotta il nuovo codec LC3, con qualità audio superiore allo SBC del vecchio standard e consumi inferiori; l'autonomia delle cuffie aumenta sensibilmente.
  \item \textbf{Audio a due canali indipendenti}: gli auricolari sinistro e destro ricevono il segnale in modo indipendente; il passaggio alla modalità mono non interrompe la connessione.
  \item \textbf{Audio broadcast Auracast}: funzione distintiva di LE Audio. Un telefono o un TV può trasmettere audio contemporaneamente a\textbf{un numero indefinito}di cuffie: nelle sale d'attesa degli aeroporti, nei musei, nelle palestre, nelle aule si può avere una sola sorgente e molte persone possono ascoltare lo stesso contenuto con le proprie cuffie.
  \item \textbf{Stato attuale}: le cuffie e i telefoni di punta usciti nel 2024-2026 supportano in genere LE Audio; i dispositivi più vecchi richiedono un aggiornamento hardware e non possono essere abilitati via software. Quando compri nuove cuffie, controlla la presenza del logo LE Audio/Auracast.
\end{itemize}

\end{tcbinfo}

\booksection{Con i sistemi per auto: un compagno sicuro alla guida}

Usare il telefono alla guida è illegale e pericoloso. Ma a volte serve davvero navigare, rispondere a una chiamata, ascoltare musica. Qui entrano in gioco CarPlay (Apple) e Android Auto (Android).

\begin{itemize}
  \item \textbf{Come collegarsi}: se la tua auto lo supporta (quasi tutte le auto nuove degli ultimi anni lo supportano), collega il telefono con un cavo dati alla porta USB dell'auto e lo schermo dell'auto si trasformerà automaticamente nell'interfaccia CarPlay o Android Auto. Alcune auto nuove supportano la connessione wireless, ancora più comoda.
  \item \textbf{Cosa puoi fare}: navigazione (molto meglio di quella integrata nell'auto), chiamate (rispondi dal volante), musica (le app musicali del telefono si proiettano sull'auto), inviare WeChat vocale (Siri o Xiao Ai leggono e rispondono ai messaggi).
  \item \textbf{Anche le auto vecchie}: se la tua auto è vecchia e non lo supporta, con qualche centinaio di yuan puoi comprare uno schermo CarPlay esterno da incollare sul cruscotto, con un'esperienza simile.
\end{itemize}

\booksection{Sicurezza e privacy nell'interazione tra dispositivi}

Più dispositivi sono collegati, più si amplia la superficie di rischio per la privacy. Alcune attenzioni:

\begin{itemize}
  \item Prima di proiettare lo schermo in pubblico, controlla che non ci siano notifiche private nella barra, per non mostrare a tutta la famiglia cose che non dovresti.
  \item Prima di regalare o vendere un dispositivo, scollegalo: vai nelle Impostazioni del telefono, trova quel dispositivo e seleziona "Rimuovi" o "Annulla abbinamento".
  \item Per le telecamere di casa intelligente, compra prodotti di marca e cifrati, non quelle a poche decine di yuan senza marca. Sono tantissime le notizie di telecamere trasmesse in diretta dagli hacker.
\end{itemize}

\begin{tcbpractice}{Prova: proietta lo schermo del telefono sul TV e crea una scena smart}
\begin{enumerate}
  \item Assicurati che telefono e TV siano sulla stessa rete Wi-Fi.
  \item Sul TV trova e attiva la funzione "Screen Mirroring" o "AirPlay".
  \item Sul telefono: iPhone tocca "Screen Mirroring" dal Centro di Controllo; Android scorri verso il basso e cerca "Proiezione wireless".
  \item Proietta sul TV una delle tue foto preferite e mostrala alla famiglia.
  \item Se hai dispositivi per la casa intelligente, prova a creare nell'app del telefono una "Modalità uscita da casa": un tocco per spegnere luci e climatizzatore e avviare il robot aspirapolvere.
\end{enumerate}

\end{tcbpractice}

\begin{tcbinfo}{Caso 1: ottimizzazione dell'autonomia della batteria in pratica}
\begin{itemize}
  \item \textbf{Scenario}: utente con Xiaomi 14, scarsa autonomia, batteria esaurita in mezza giornata.
  \item \textbf{Passaggi}: (1) Impostazioni Batteria classifica dei consumi (2) identifica app anomale in background (3) disattiva 5G e passa a 4G (4) limita il background (5) attiva la modalità risparmio energetico (6) durante la ricarica togli la cover per dissipare il calore.
  \item \textbf{Risultato atteso}: autonomia prolungata del 30-50\%, dato di stato di salute della batteria leggibile.
\end{itemize}

\end{tcbinfo}
